
\subsection{Теоретические основы формирования и анализа морских волн}
\label{sec:wave_theory}

Естественным началом к описанию анализа влияния морских волн на берег является описание механизма возникновения и развития волн, рассмотренного в работе~\cite{DeanDalrymple1991}.

Неоднородный прогрев атмосферы создает перепады атмосферного давления, формируя потоки ветра.
Проходя над водной поверхностью ветер, за счет сил трения начинает формировать сдвиговые напряжения, возмущающие гладкую поверхность моря и порождая рябь -- короткие капиллярные волны.
С появлением ряби шероховатость поверхности возрастает, что позволяет ветру эффективнее передавать воде импульс, усиливая волны.
Развивающиеся волны поднимаются выше и принимают еще больший импульс, усиливаясь еще сильнее, в результате формируя собственный импульс и начиная движение, согласованное с направлением ветра.
Усиление и рост волн происходит тем сильнее, чем выше скорость ветра, чем дольше он дует и чем больше расстояние согласованого с волнами движения ветра над водой (\textit{fetch} -- непрерывной трасса ветра над морем).

Влияние порождающих волны факторов конечно и приводит в итоге к \enquote{развитому волнению} (Fully Developed Sea (FDS)), на стадии которого система волн становится стохастической.
На этом этапе происходят процессы интерференции волн, слипания (разрушения) их гребней и перераспределение энергии по волновому спектру из-за чего часть волн поднимается, становясь при этом более длинными.
Если ветер ослабевает или меняет направление, остается зыбь, а более длинные, ровные волны, пришедшие из области штормов и распространяются практически без потерь на большие расстояния, доходя до берега.

Проходя над глубоководьем, частицы воды в волне совершают почти замкнутые по орбитам движения и форма волны подвержена малым изменениям, но по мере приближения волны к берегу начинается этап ее трансформации.
По мере того, как волна заходит на мелководье, нижняя часть ее орбиты начинает пересекать поверхность дна, тормозясь при этом.
На этом этапе запускаются два процесса трансформации волны на мелководье.
На первом, длина волны сокращается, что приводит к росту высоты, делая фронт волны круче.
На втором, за счет неравномерного торможения о дно волна начинает выравниваться по направлению к берегу.
В результате, достигнув предельной крутизны, волна теряет устойчивость: ее гребень перекидывается вперед и происходит обрушение, то, что мы наблюдаем как прибой в прибрежной зоне~\cite{DeanDalrymple1991}.

Опираясь на описанный механизм возникновения волн, становится возможным выделить следующие этапы расчета:
\begin{itemize}
    \item расчет параметров волны в глубокой воде;
    \item корректировка параметров волн за счет трансформации на мелководье;
    \item учет изменения направления волны за счет рефракции;
    \item расчет удельной мощности волны приходящейся на берег.
\end{itemize}

Разберем перечисленные этапы подробнее.



% \subsubsection*{Этап подготовки данных}

% На первом этапе необходимо выполнить сбор и подготовку исходных данных.
% Минимально достаточными исходными данными для анализа волновых нагрузок будут являться:
% \begin{itemize}
% \item исторические метедоанные, содержащие информацию о скорости и направлении ветра;
% \item данные о общем положении береговой линии, необходиой для расчета длины разгона волн (fetch);
% \item данные о батиметрии в области анализируемой точки.
% \end{itemize}

% Источники и способы получения исходных данных, а так же их полная спецификация будет рассмотрена в соответствующих разделах (№№ - №№) далее.
% Однако уже на этом этапе необходимо определить к ним базовые требования.
% Так как на величину параметров волн в открытой, глубокой воле непосредственно влияет продолжительность постоянного ветра максимальная точность результатов достигается за счет интегрирования результатов максимально коротких временных интервалов между метеоданными.
% Опыт расчета по заранее агрегированным розам ветров показал свою недостаточную точность, вследствие чего в метод расчета заложено использование суточных методанных данных.

% Так как для вычисления длинны разгона волн необходимо знание положения береговой во всех направлениях относительно точки расчета исходной трассы береговой линии, задаваемой пользователем на этапе формулирования задания, недостаточно, в следствие чего необходим сбор дополнительных. глобальных данных о морском побережье.

% Участвующая в корректировке параметров волн батиметрия необходима для определения границы "мелководья" вдоль рассчтываемого направления ветра.
% Технически этап сбора и использования батиметрических данных может быть заменен вводом условного "радиуса" по дуге которого омжно условно провести границу мелководья.
% Несмотря на то, что подобное допущение сокращает объем собираемой информации и минимизирует объем выполняемых вычислений опыт выполненных расчетов показал, что экономия затрачиваемых ресурсов не оправдывает необходимость априорного задания от пользователя дополнительного гиперпараметра. 

% \subsubsection{Вычисление длинны разгона волн}
